\documentclass[11pt]{article}

% === Packages ===
\usepackage[utf8]{inputenc}
\usepackage[T1]{fontenc}
\usepackage{geometry}
\geometry{a4paper, margin=1in}
\usepackage{graphicx}
\usepackage{booktabs}
\usepackage{amsmath}
\usepackage{amsfonts}
\usepackage{hyperref}
\usepackage{xcolor}
\usepackage{enumitem}
\usepackage{float}

\hypersetup{
    colorlinks=true,
    linkcolor=blue,
    urlcolor=blue,
    citecolor=blue,
}

% === Title ===
\title{\textbf{How Far Can Aerial Detection Go on 8 GB?\\A Systematic Ablation Study of Aerial Small-Object Detection}}

\author{
    Haoyi (Yuki) Zou\\[2mm]
    \small Independent Research\\[1mm]
    \small \texttt{https://github.com/ZHY9981/aerial-yolo26s}
}

\date{July 2026}

\begin{document}
\maketitle

\begin{abstract}
Aerial small-object detection poses a unique challenge: targets smaller than 30 pixels must be
detected in high-resolution drone footage, while consumer GPUs (8 GB VRAM) cannot accommodate
the high-resolution detection heads that small objects require. We present a customized YOLO26s
architecture that addresses three compounding bottlenecks—positional information loss, small-object
gradient starvation, and VRAM overcommitment—through a combination of Coordinate Attention (CoordAtt),
a P3+P4 dual-head design, WIoU v3 loss, and a reduced TaskAlignedAssigner pixel threshold. Across
20+ systematically ablated experiments on a curated 7-class aerial dataset, our best model (V16.0)
achieves 74.00\% mAP@0.5 and 52.11\% mAP@0.5:0.95 with only 7.03M parameters, operating within
an 8 GB VRAM budget at 4.0 ms inference latency. We also document four failed experiments and six
key insights, including the counterintuitive finding that a dual-head architecture outperforms
triple-head by 0.71\% mAP under memory constraints, in a controlled comparison.
\end{abstract}

% ======================================================================
\section{Introduction}
% ======================================================================

Aerial object detection from drone-captured imagery is a problem of growing practical importance,
with applications ranging from traffic surveillance to search-and-rescue operations. However, it
exposes a fundamental tension in modern object detectors: small targets require high-resolution
feature maps, but high-resolution feature maps require GPU memory that consumer hardware does
not have.

Consider a pedestrian viewed from 100 meters altitude. At 800-pixel input resolution, this target
occupies approximately 20$\times$15 pixels—less than 0.05\% of the image area. A standard YOLO
detector with P3/8, P4/16, and P5/32 downsampling strides places this target on the P3 feature map
(100$\times$100 spatial resolution), where each cell corresponds to an 8$\times$8 pixel region in
the original image. A 20-pixel target spans roughly 2.5 cells—barely enough for the detector to
localize it. Adding a P2/4 head (200$\times$200 resolution) would help, but the TaskAlignedAssigner's
alignment matrix at P2 resolution exceeds 8 GB VRAM at batch size 4.

This paper documents a systematic attempt to maximize aerial detection performance under a fixed
hardware constraint (RTX 5060 Laptop GPU, 8 GB VRAM). Rather than scaling hardware, every design
decision is evaluated by both accuracy and memory cost. The key contributions are:

\begin{enumerate}[leftmargin=*]
    \item A per-scale Coordinate Attention (CoordAtt) injection strategy that preserves fine-grained
          spatial cues for small targets without the VRAM cost of a full P2 head.
    \item A P3+P4 dual-head architecture that outperforms both the standard P3+P4+P5 triple-head
          and the P2+P3+P4+P5 quad-head on 8 GB hardware.
    \item A 20-version ablation study isolating the effect of each architectural modification,
          loss function, and data preprocessing decision with all failure cases documented.
\end{enumerate}

% ======================================================================
\section{Problem Analysis}
% ======================================================================

Aerial small-object detection faces three compounding difficulties:

\subsection{Extreme Scale Variation}

In aerial imagery, a freight truck (occupying 300$\times$100 pixels at 50 m altitude) and a
pedestrian (20$\times$15 pixels) coexist in the same image. The detector must simultaneously
handle targets spanning two orders of magnitude in pixel area. Standard multi-scale detection
heads (P3 through P5) are designed for COCO-scale objects, where the smallest targets (e.g.,
tennis rackets, bottles) are still $\sim$50 pixels—double the size of an aerial pedestrian.

\subsection{Memory Constraints on Consumer Hardware}

The RTX 5060 Laptop GPU provides 8 GB of VRAM. At batch size 4 and image size 800$\times$800,
the standard YOLO26s triple-head configuration consumes approximately 6.8 GB during training,
leaving only 1.2 GB of headroom. Adding a P2 head (feature map 200$\times$200$\times$25,600 cells)
causes the TaskAlignedAssigner to allocate a $B \times N_{\text{anchors}} \times N_{\text{gt}}$
alignment matrix: at P2 resolution with 7 classes and up to 200 ground-truth instances per image,
this matrix alone exceeds 2.5 GB—pushing total VRAM usage past the 8 GB limit and causing an OOM
error.

\subsection{Class Imbalance}

The aerial\_v9 dataset exhibits extreme long-tail distribution: person (69,061 instances) and
car (68,280) dominate, while freight (1,180) and small-bus (987) each represent less than 2\%
of training instances. Standard equal-weight classification loss causes the model to optimize
almost exclusively for the majority classes.

% ======================================================================
\section{Methodology}
% ======================================================================

\subsection{Coordinate Attention (CoordAtt)}

Channel attention mechanisms like Squeeze-and-Excitation (SE) improve feature representation by
modeling inter-channel dependencies, but they discard all spatial information through global
average pooling. For small aerial targets, \textit{where} a feature is located is as important
as \textit{what} it represents.

CoordAtt \cite{hou2021coordinate} decomposes 2D global pooling into two 1D feature encoding
operations along the horizontal and vertical directions:

\begin{align}
    \mathbf{z}^h_c(h) &= \frac{1}{W} \sum_{0 \leq i < W} \mathbf{x}_c(h, i), \quad
    \mathbf{z}^w_c(w) = \frac{1}{H} \sum_{0 \leq j < H} \mathbf{x}_c(j, w)
\end{align}

The two 1D encodings are concatenated, passed through a shared 1$\times$1 convolution for
cross-channel interaction, then split into separate attention maps $\mathbf{a}^h$ and
$\mathbf{a}^w$. The attended output is:

\begin{align}
    \mathbf{y}_c(i, j) = \mathbf{x}_c(i, j) \times \mathbf{a}^h_c(i) \times \mathbf{a}^w_c(j)
\end{align}

\textbf{Placement strategy.} We place CoordAtt per-scale at P3 and P4 neck positions (immediately
before the Detect head), rather than at the backbone P5 layer. Our V16 ablation (Section~\ref{sec:experiments})
shows this is critical: a backbone P5 CoordAtt must propagate signals through two nearest-neighbor
upsample steps to reach P3 detection features, losing spatial precision. Per-scale placement
directly enhances the 120$\times$120 (P3) and 60$\times$60 (P4) feature maps.

\subsection{P3+P4 Dual-Head Architecture}

Removing the P5 detection head reduces the parameter count from 8.06\,M (triple-head, V20)
to 7.03\,M (dual-head, V16) and frees a portion of detection-head VRAM. This freed memory is
reallocated to larger P3/P4 feature channels, improving small-object feature quality.

The counterintuitive result is that the dual-head configuration \textit{improves} mAP over
triple-head. In our controlled comparison — V16 (dual-head) versus V20 (triple-head), both with
per-scale CoordAtt at imgsz=800, differing \textit{only} in head count — the dual-head variant
gains \textbf{+0.71\%} mAP@0.5. This is not because P5 is useless (the P5 head detects large
objects such as buses and freight trucks effectively), but because under the 8 GB memory
constraint, the memory saved by removing P5 allows larger P3/P4 channels that improve detection
of the majority of targets (person, cycle, car, truck), which collectively account for over
80\% of instances.

\begin{quote}
\textbf{Note on effect size.} An earlier version pair (V12 $\rightarrow$ V14) exhibits a larger
+2.1\% gap, but that comparison changed five variables simultaneously (head count, TAL threshold,
class weighting, augmentation scale, and imgsz). We report the controlled V16-vs-V20 result
(+0.71\%) as the honest isolated effect of the head-count change.
\end{quote}

\subsection{WIoU v3 Loss}

Standard CIoU loss treats all bounding box regression errors equally. WIoU v3 \cite{tong2023wise}
introduces an attention-based outlier suppression mechanism that adaptively weights regression
gradients based on the quality of each prediction:

\begin{align}
    \mathcal{L}_{\text{WIoU v3}} = r \cdot \mathcal{L}_{\text{IoU}},
    \quad r = \frac{\beta}{\delta \alpha^{\beta - \delta}}
\end{align}

where $\beta = \mathcal{L}_{\text{IoU}}^* / \overline{\mathcal{L}_{\text{IoU}}}$ measures
outlier degree, and $\alpha, \delta$ are hyperparameters controlling suppression strength.
Small aerial targets (high $\beta$ due to difficult localization) receive amplified gradients,
while large targets (low $\beta$) are unaffected. This modification costs zero additional VRAM.

\subsection{TAL Small-GT Threshold}

The TaskAlignedAssigner (TAL) \cite{feng2021toad} can fail on ground-truth boxes smaller than
the smallest feature stride: such boxes may contain no anchor center at all and therefore
receive no positive assignment. TAL handles this by expanding the affected boxes to
\texttt{stride\_val} so at least one anchor lands inside. The default cutoff for this rescue
is \texttt{stride[0]} (8 px at P3).

We halve this cutoff to \texttt{stride[0]/2} (4 px), extending the rescue to boxes in the
\textbf{4--8 px} range — the size band where the smallest aerial targets (pedestrians, cyclists)
concentrate. This yields approximately 4$\times$ more positive anchors for the smallest targets
and improves person/cycle recall by +6--7\% each.

\subsection{Class-Weighted Loss}

We apply per-class multiplicative weights to the classification loss. The weight vector is
indexed by the class order defined in \texttt{data.yaml}:

\begin{equation}
    \mathbf{w} = [2.0, 3.0, 1.8, 1.5, 1.0, 1.0, 1.0]
\end{equation}

corresponding to \texttt{[person, cycle, bus, small-bus, car, truck, freight]} — that is,
the effective weights are person~$\times$2.0, cycle~$\times$3.0, bus~$\times$1.8,
small-bus~$\times$1.5, with car/truck/freight at base weight. Cycle receives the highest weight
(3.0$\times$) because it is both rare (14K instances) and small; person is up-weighted as the
hardest small-object class.

% ======================================================================
\section{Experiments}
\label{sec:experiments}
% ======================================================================

\subsection{Dataset}

The aerial\_v9 dataset is curated from three public sources: aerial.v1i (Roboflow, CC BY 4.0),
Aerial Vehicle Detection (MIT), and aerial.v3i (MIT). The curation process involved merging the
three sources, unifying class taxonomies (e.g., motorcycle + bicycle $\rightarrow$ cycle,
van $\rightarrow$ car), and manually inspecting approximately 500 random samples to remove label
noise and add missing annotations. The final dataset contains 8,075 training, 2,224 validation,
and 738 test images across 7 classes.

\subsection{Training Configuration}

All experiments use a fixed hardware and training protocol:

\begin{table}[H]
\centering
\caption{Training configuration (all versions).}
\begin{tabular}{ll}
\toprule
GPU & RTX 5060 Laptop 8 GB \\
Batch size & 4 (P2 experiments: 2) \\
Optimizer & SGD, lr$_{0}$=0.01, momentum=0.937, weight\_decay=0.0005 \\
Epochs & 200 (V8--V16), 120 (V17--V20) \\
Image size & 640--960 (V16 uses 800) \\
Augmentation & mosaic=0.5, scale=0.3, no mixup/copy-paste \\
Loss & WIoU v3 (box), VFL (cls), DFL \\
\bottomrule
\end{tabular}
\end{table}

Mixup and copy-paste augmentation are deliberately disabled for aerial data—our early experiments
(V1--V3) showed they degrade performance by 1.5--2.0\%, likely because they create unrealistic
object placements (e.g., a car floating on a rooftop) when applied to aerial imagery with dense,
structured layouts.

\subsection{Ablation Study}

Table~\ref{tab:ablation} presents the key milestones from the 20-version ablation study. Each
version modifies exactly one variable from the preceding baseline.

\begin{table}[H]
\centering
\caption{Key ablation results. Full table in project documentation.}
\label{tab:ablation}
\begin{tabular}{clccc}
\toprule
Version & Key Change & mAP@0.5 & mAP@0.5:0.95 & $\Delta$ \\
\midrule
V8.0  & Baseline (clean data) & 67.81\% & 44.31\% & --- \\
V8.1  & + WIoU v3 loss & 69.59\% & 47.00\% & +1.78\% \\
V10.0 & + CoordAtt + aerial\_v9 & 70.51\% & 49.33\% & +0.92\% \\
V12.0 & imgsz 640$\rightarrow$800 & 71.58\% & 50.27\% & +1.07\% \\
V14.0 & P3+P4 dual-head, imgsz=960 & 73.71\% & \textbf{52.65\%} & +2.13\% \\
\textbf{V16.0} & \textbf{P3+P4 + per-scale CoordAtt, 800} & \textbf{74.00\%} & 52.11\% & +0.29\% \\
V17.0 & + P4 RepNCSPELAN4 & 73.52\% & 51.27\% & $-$0.48\% \\
V18.0 & + ECA / ASFF & 72.91\% & 50.84\% & $-$1.09\% \\
V19.0 & P4 wide channel 256$\rightarrow$384 & 73.44\% & 51.63\% & $-$0.56\% \\
V20.0 & P3+P4+P5 triple-head (clean) & 73.29\% & 50.41\% & $-$0.71\% \\
\bottomrule
\end{tabular}
\end{table}

Table~\ref{tab:perclass} shows the per-class breakdown of V16.0, our production model.

\begin{table}[H]
\centering
\caption{Per-class performance — V16.0 (best model).}
\label{tab:perclass}
\begin{tabular}{lcccc}
\toprule
Class & AP50 & AP50-95 & Precision & Recall \\
\midrule
person    & 60.2\% & 24.7\% & 77.7\% & 50.9\% \\
cycle     & 43.0\% & 17.6\% & 67.0\% & 36.3\% \\
bus       & 91.8\% & 74.5\% & 91.3\% & 86.3\% \\
small-bus & 99.2\% & 83.3\% & 92.3\% & 99.4\% \\
car       & 67.2\% & 47.3\% & 87.2\% & 51.4\% \\
truck     & 60.0\% & 41.0\% & 63.8\% & 53.6\% \\
freight   & 96.6\% & 76.6\% & 89.8\% & 95.0\% \\
\midrule
\textbf{Overall} & \textbf{74.0\%} & \textbf{52.1\%} & \textbf{81.3\%} & \textbf{67.6\%} \\
\bottomrule
\end{tabular}
\end{table}

\subsection{Failed Experiments}

Four experiments failed to improve performance, and we document them as honest records:

\begin{itemize}[leftmargin=*]
    \item \textbf{V4.0 -- Freeze+unfreeze 2-stage training (46.82\% mAP@0.5).}
          Standard YOLO transfer-learning protocol (freeze backbone, train head, unfreeze all)
          performs worse than training from scratch. Hypothesis: COCO pretrained features encode
          object shape priors (e.g., ``person = upright rectangle'') that are orthogonal to the
          overhead view of aerial imagery, where pedestrians appear as small blobs.
    \item \textbf{V7.0 -- VisDrone noisy labels (55.25\% mAP@0.5, Recall 50\%).}
          VisDrone2019 annotations contain significant label noise (motorcycles mislabeled as
          bicycles, inconsistent labeling across frames). The 12.56\% gap between V7 and V8
          demonstrates that data quality trumps architectural sophistication.
    \item \textbf{V8.2 -- Transfer-learning relay (69.71\% mAP@0.5).}
          Fine-tuning V8.1 weights on the expanded aerial\_v9 dataset with the default learning
          rate caused catastrophic forgetting—all 7 classes degraded simultaneously. The learning
          rate reset on a larger dataset destroyed learned feature representations.
    \item \textbf{V13.0 -- P2 high-resolution head (66.68\% mAP@0.5, OOM).}
          While theoretically beneficial for small objects, the P2 head's 25,600-cell feature
          map exceeds 8 GB VRAM with the TaskAlignedAssigner. The OOM occurs during assignment
          matrix allocation, before any forward pass.
\end{itemize}

% ======================================================================
\section{Key Findings}
% ======================================================================

\begin{enumerate}[leftmargin=*]
    \item \textbf{Data quality beats architecture.} Switching from noisy VisDrone to curated
          aerial\_v9 yielded +12.56\% mAP@0.5 (V7 $\rightarrow$ V8)—more than all four
          architectural modifications combined.
    \item \textbf{Dual-head outperforms triple-head on 8 GB.} In a controlled comparison (V16
          vs V20, only head count differs), the dual-head variant gains +0.71\% mAP while using
          about 13\% fewer parameters. The memory saved enables larger P3/P4 channels that
          benefit the most common target classes.
    \item \textbf{WIoU v3 is the highest-ROI change.} +1.78\% mAP@0.5 with zero VRAM cost and
          approximately 5 lines of code changed.
    \item \textbf{CoordAtt is cheap but effective.} Negligible parameter increase (+0.2M) for
          +2.4\% Precision gain.
    \item \textbf{TAL threshold directly boosts small targets.} A simple 8$\rightarrow$4 pixel
          threshold change improves person and cycle by +6--7\% each.
    \item \textbf{Mixup/copy-paste harm aerial detection.} These augmentations degrade performance
          by 1.5--2.0\% on aerial data, likely due to unrealistic spatial configurations.
\end{enumerate}

% ======================================================================
\section{Limitations and Future Work}
% ======================================================================

\subsection{Car Detection Bottleneck}

Car detection on V16.0 (67.3\% AP50) lags person (60.2\%) in absolute terms but represents the
largest opportunity for improvement given car's high instance count (68,280). Current hypotheses:

\begin{itemize}[leftmargin=*]
    \item \textbf{Scale ambiguity:} Cars viewed from 50--100 m altitude span a wide range of
          apparent sizes (50--200 pixels) that may confuse the P4-level detector. A P2 head
          would help differentiate small cars from large cars, but remains OOM-bound.
    \item \textbf{Occlusion patterns:} Cars in urban aerial scenes are frequently partially
          occluded by trees, buildings, or other vehicles. The current augmentation pipeline
          does not include occlusion simulation.
\end{itemize}

Ongoing work explores SimOTA center-prior assignment and SGLoss-style adaptive grid selection
\cite{li2022generalized} to recover P2-head benefits without the OOM cost.

\subsection{Rare Class Generalization}

Freight and small-bus each represent <2\% of training instances. While their detection
performance is high on the current validation set (96.6\% and 99.2\% AP50 respectively),
these metrics reflect the limited visual diversity of rare-class examples. Cross-dataset
generalization remains untested.

\subsection{Field-Deployment Constraints}

The current pipeline achieves 4.0 ms inference on RTX 5060 (batch=1, FP16, 800$\times$800),
which is sufficient for real-time analysis on a consumer laptop in an outdoor setting. The
motivating use case — running detection on a single laptop GPU without cloud or desktop
hardware — is met by this configuration. TensorRT quantization is planned for further
latency reduction on lower-power edge devices.

% ======================================================================
\section{Conclusion}
% ======================================================================

This project demonstrates that consumer-grade hardware (RTX 5060 8 GB) is sufficient for
competitive aerial small-object detection when architectural decisions are co-designed with
memory constraints. The key insight is not that any single modification is revolutionary, but
that four targeted changes—CoordAtt, dual-head design, WIoU v3, and TAL threshold tuning—each
address a specific bottleneck ($\sim$2\% improvement) and compound to a meaningful overall gain
(+6.19\% from V8.0 baseline). The 20-version ablation log, including documented failures,
provides a reproducible baseline for future work on resource-constrained aerial detection.

% ======================================================================
% References
% ======================================================================

\begin{thebibliography}{9}

\bibitem{hou2021coordinate}
Q.~Hou, D.~Zhou, and J.~Feng,
``Coordinate Attention for Efficient Mobile Network Design,''
in \emph{Proc. CVPR}, 2021, pp. 13713--13722.

\bibitem{tong2023wise}
Z.~Tong, Y.~Chen, Z.~Xu, and R.~Yu,
``Wise-IoU: Bounding Box Regression Loss with Dyna
\bibitem{feng2021toad}
C.~Feng, Y.~Zhong, Y.~Gao, M.~R.~Scott, and W.~Huang,
``TOOD: Task-aligned One-stage Object Detection,
in \emph{Proc. ICCV}, 2021, pp. 3490--3499.

\bibitem{li2022generalized}
X.~Li, W.~Wang, L.~Wu, S.~Chen, X.~Hu, J.~Li, J.~Tang, and J.~Yang,
``Generalized Focal Loss: Learning Qualified and Distributed Bounding Boxes for Dense Object Detection,
in \emph{Proc. NeurIPS}, 2020.

\bibitem{jocher2023ultralytics}
G.~Jocher, A.~Chaurasia, and J.~Qiu,
``Ultralytics YOLO,
\emph{Ultralytics}, 2023. [Online]. Available: \url{https://github.com/ultralytics/ultralytics}

\end{thebibliography}

\end{document}
